../cictikz/src/cictikz/data/tex/cictikz_preamble.tex

% Pierce crystal oscillator: an inverting transconductance (-gm) with a
% feedback resistor R1, a series resistor R2 on the output, the crystal
% between the two pins, and a load capacitor on each side. The dotted
% line is the chip boundary: the inverter, R1 and R2 are on the IC, the
% crystal and the load capacitors are on the PCB.

\begin{circuitikz}[american, thick, transform shape, circuit ee IEC]

  ../cictikz/src/cictikz/data/tex/cictikz_lib.tex

  \def\xl{0}     % left (input) node
  \def\xr{4.4}   % right (output) node
  \def\yinv{2.9} % inverter row
  \def\yr{4.6}   % feedback resistor row
  \def\yx{0}     % crystal row
  \def\yc{-1.0}  % top of load capacitors

  % inverting transconductance
  \node[not port] (inv) at ({\xl/2+\xr/2},\yinv) {};
  \draw (\xl,\yinv) -- (inv.in);
  \draw (inv.out) -- (\xr,\yinv);
  \node[anchor=north] at ({\xl/2+\xr/2},\yinv-0.65) {$-g_m$};

  % feedback resistor
  \draw (\xl,\yr) to [R, l=$R_1$] (\xr,\yr);
  \draw (\xl,\yr) -- (\xl,\yx);
  \fill (\xl,\yinv) circle (0.075);

  % series resistor on the output
  \draw (\xr,\yr) -- (\xr,\yinv);
  \fill (\xr,\yinv) circle (0.075);
  \draw (\xr,\yinv) to [R, l=$R_2$] (\xr,1.2) -- (\xr,\yx);

  % crystal
  \draw (\xl,\yx) -- (1.5,\yx)
        (1.5,\yx-0.35) -- (1.5,\yx+0.35)
        (1.75,\yx-0.5) rectangle (2.65,\yx+0.5)
        (2.9,\yx-0.35) -- (2.9,\yx+0.35)
        (2.9,\yx) -- (\xr,\yx);

  % load capacitors
  \draw (\xl,\yx) to [short,*-] (\xl,\yc);
  \draw (\xl,-2.6) \vground \vcapacitor{};
  \draw (\xr,\yx) to [short,*-] (\xr,\yc);
  \draw (\xr,-2.6) \vground \vcapacitor{};

  % chip boundary
  \draw[densely dotted] (-1.2,1.5) .. controls (1.2,1.5) and (2.2,0.7)
    .. (5.9,0.7);
  \node[anchor=south west] at (5.3,0.8) {IC};
  \node[anchor=north west] at (5.3,0.6) {PCB};

\end{circuitikz}
\end{document}
